% UBT-AI-PROVENANCE-BEGIN
% schema: ubt-ai-provenance/v1
% tier: C_working
% ai_assistance: disclosed
% human_review: risk-based
% editorial_responsibility: Ing. David Jaroš
% policy: ../../../AI_PROVENANCE.md
% notice: Working material; exhaustive human review is not claimed.
% UBT-AI-PROVENANCE-END

\chapter{Testy, reprodukovatelnost a formální ověřování}
\label{ch:tests-repro}

Živou ověřovací vrstvou je aktivní repozitář, nikoli \texttt{ARCHIVE/}.
Důležité vstupní body jsou:
\begin{itemize}
 \item \texttt{tests/} pro vědecké, provenienční a architektonické regresní testy;
 \item \texttt{tools/apply\_provenance\_headers.py} pro značení AI provenance;
 \item \texttt{tools/latex\_audit.py} pro izolované překlady LaTeX dokumentů;
 \item \texttt{tools/check\_alpha\_claims.py} a aktivní testy programu alfa pro
       kontrolu tvrzení citlivých na aktuální stav;
 \item numerické kontroly větve theta/Mellin/Gauss zajišťuje\\
       \texttt{research\_tracks/theta\_spectral/}\\
       \texttt{theta\_zeta\_probe.py};
 \item \texttt{docs/DERIVATION\_VERIFICATION\_POLICY.md} jako závazná pravidla
       kontroly odvození v článcích.
\end{itemize}

\section{Čtyři různé druhy kontroly}
Je důležité nezaměňovat různé úrovně ověření. Překlad LaTeXu pouze dokazuje,
že je dokument syntakticky sestavitelný. Numerický test ověřuje konkrétní
číselné důsledky v dané přesnosti. CAS nástroj, například SymPy nebo Maxima,
ověřuje symbolickou identitu, kterou jsme mu skutečně zadali. Formální systém
Lean kontroluje, zda formalizovaný závěr plyne z formalizovaných předpokladů.
Ani jedna z těchto kontrol sama o sobě nedokazuje, že příslušné předpoklady jsou
fyzikálně vybrány kanonickou UBT dynamikou.

\section{Lean jako preferovaný formální kontrolor}
Pro nová nebo podstatně změněná teorémová odvození je preferovaným nástrojem
Lean. Pokud je tvrzení rozumně formalizovatelné, cílem je mít vedle analytického
odvození také kompilující \texttt{.lean} důkaz. Jestliže formalizace ještě není
hotová nebo v daném prostředí Lean není dostupný, musí být stav výslovně
označen jako \texttt{LEAN-PENDING}; nesmí se napsat, že tvrzení je ``formálně
dokázáno'' jen proto, že agent vygeneroval Lean syntaxi.

Lean nenahrazuje ostatní kontroly. U hlavního výsledku článku je žádoucí druhý,
nezávislý kanál: například SymPy/Maxima pro symbolickou algebru a
NumPy/Octave/MATLAB/Mathcad pro numerickou reprodukci. Dvě implementace by měly
pokud možno používat odlišnou formulaci, aby pouze neopakovaly stejnou chybu.

\section{Povinný záznam ověření článku}
Každý nový nebo podstatně změněný článek má obsahovat oddíl \emph{Ověření}
nebo odkaz na doprovodný záznam. Pro hlavní tvrzení se zapisuje alespoň:
\begin{itemize}
 \item identifikátor tvrzení, rovnice nebo lemmatu;
 \item analytický stav a předpoklady;
 \item použitý nástroj a jeho verze;
 \item cesta k reprodukčnímu skriptu, testu, pracovnímu listu nebo \texttt{.lean}
       souboru;
 \item výsledek a případná tolerance;
 \item přesně co kontrola ověřuje a co naopak neověřuje;
 \item stav Lean formalizace:\texttt{PROVED},\texttt{PARTIAL},
       \texttt{LEAN-PENDING} nebo \texttt{NOT-APPLICABLE}.
\end{itemize}

\section{Průběžná integrace učebnice}
Vyhrazený pracovní postup je \texttt{.github/workflows/textbook.yml}. Spouští
se při relevantních odesláních změn a požadavcích na sloučení a lze jej spustit
také ručně. Nejprve
kontroluje shodu jazykových verzí a poté používá
\texttt{tools/latex\_audit.py --strict}. Chyba kteréhokoli samostatného
dokumentu učebnice proto způsobí selhání pracovního postupu.

Po přísném auditu pracovní postup volá stejný Makefile jako při lokálním sestavení a
vytvoří přesně dvě oficiální PDF v \texttt{build/textbook/public/}:
\begin{itemize}
 \item \texttt{UBT\_Textbook\_EN.pdf} --- anglická edice;
 \item \texttt{UBT\_Textbook\_CS.pdf} --- česká edice.
\end{itemize}
Každé PDF se nahraje jako samostatně pojmenovaný GitHub Actions artefakt.
E-mailová distribuce přikládá obě edice a je povolena pouze po úspěšném
odeslání změn do \texttt{master}; požadavek na sloučení ani ruční spuštění
e-mail nikdy neodesílají.

Pro lokální publikovanou kopii lze spustit
\begin{verbatim}
make -C docs/textbook publish
\end{verbatim}
což zkopíruje stejnou dvojici do \texttt{docs/pdfs/}. Vygenerovaná PDF jsou
pouze publikační výstupy; vědecký stav zůstává řízen zdrojovým textem a
registrem provenience a stavu.

\section{Minimální lokální kontrola před textbook patchem}
Doporučená sekvence je:
\begin{verbatim}
python tools/apply_provenance_headers.py --apply
python tools/apply_provenance_headers.py --check
python tools/latex_audit.py --include 'docs/textbook/*.tex' \
  --include 'docs/textbook/**/*.tex' --jobs 2 --timeout 240 --strict
make -C docs/textbook verify
\end{verbatim}
Úspěšné sestavení potvrzuje syntaktické a reprodukční vlastnosti. Nezvyšuje samo od
sebe tvrzení úrovně C na dokázaný výsledek UBT.
